\documentclass[a4paper,12pt]{article}
\usepackage[utf8]{inputenc}
\usepackage[russian]{babel}
\usepackage{amsmath, amssymb}
\usepackage{geometry}
\geometry{a4paper, left=25mm, right=25mm, top=20mm, bottom=20mm}

\begin{document}

\begin{center}
    \Large\textbf{Отчет по лабораторной работе} \\
    \large\textbf{«Стабилизация перевернутого маятника»}
\end{center}

\section*{1. Цель работы}
Исследование возможности стабилизации перевернутого маятника в верхнем положении с помощью горизонтального перемещения точки опоры. Анализ различных стратегий управления.

\section*{2. Описание системы}
Рассматривается математический маятник:
\begin{itemize}
    \item масса $m$,
    \item длина нити $l$,
    \item точка подвеса может перемещаться горизонтально.
\end{itemize}
Введем координаты:
\begin{itemize}
    \item $\varphi$ — угол отклонения маятника от вертикали,
    \item $u$ — горизонтальное смещение точки опоры.
\end{itemize}
Управление: ускорение точки опоры $\ddot{u}$.

\section*{3. Вывод уравнений движения}

\subsection*{3.1. Координаты маятника}
\[
x = u + l\sin\varphi, \qquad y = l\cos\varphi.
\]

\subsection*{3.2. Скорости}
\[
\dot{x} = \dot{u} + l\dot{\varphi}\cos\varphi, \qquad \dot{y} = -l\dot{\varphi}\sin\varphi.
\]

\subsection*{3.3. Кинетическая энергия}
\[
T = \frac{m}{2}(\dot{x}^2 + \dot{y}^2) = \frac{m}{2}\left(\dot{u}^2 + 2l\dot{u}\dot{\varphi}\cos\varphi + l^2\dot{\varphi}^2\right).
\]

\subsection*{3.4. Потенциальная энергия}
Нулевой уровень — точка подвеса:
\[
\Pi = mgy = mgl\cos\varphi.
\]

\subsection*{3.5. Функция Лагранжа}
\[
L = T - \Pi = \frac{m}{2}\left(\dot{u}^2 + 2l\dot{u}\dot{\varphi}\cos\varphi + l^2\dot{\varphi}^2\right) - mgl\cos\varphi.
\]

\subsection*{3.6. Уравнение Лагранжа для $\varphi$}
\[
\frac{d}{dt}\left(\frac{\partial L}{\partial \dot{\varphi}}\right) - \frac{\partial L}{\partial \varphi} = 0.
\]
Вычисляем:
\[
\frac{\partial L}{\partial \dot{\varphi}} = ml\dot{u}\cos\varphi + ml^2\dot{\varphi},
\]
\[
\frac{d}{dt}\left(\frac{\partial L}{\partial \dot{\varphi}}\right) = ml\ddot{u}\cos\varphi - ml\dot{u}\dot{\varphi}\sin\varphi + ml^2\ddot{\varphi},
\]
\[
\frac{\partial L}{\partial \varphi} = -ml\dot{u}\dot{\varphi}\sin\varphi + mgl\sin\varphi.
\]
После подстановки и сокращения:
\[
ml\ddot{u}\cos\varphi + ml^2\ddot{\varphi} - mgl\sin\varphi = 0.
\]
Делим на $ml$:
\[
\ddot{u}\cos\varphi + l\ddot{\varphi} - g\sin\varphi = 0.
\]
Окончательно:
\[
l\ddot{\varphi} + \ddot{u}\cos\varphi - g\sin\varphi = 0,
\]
или
\[
\ddot{\varphi} - \frac{g}{l}\sin\varphi = -\frac{\ddot{u}}{l}\cos\varphi. \tag{1}
\]

\subsection*{3.7. Линеаризация}
При малых $\varphi$: $\sin\varphi \approx \varphi$, $\cos\varphi \approx 1$. Получаем:
\[
\ddot{\varphi} - \frac{g}{l}\varphi = -\frac{\ddot{u}}{l}. \tag{2}
\]

\section*{4. Неуправляемый маятник}
При $\ddot{u} = 0$:
\[
\ddot{\varphi} - \frac{g}{l}\varphi = 0.
\]
Характеристическое уравнение:
\[
\lambda^2 - \frac{g}{l} = 0 \quad\Rightarrow\quad \lambda_{1,2} = \pm\sqrt{\frac{g}{l}}.
\]
Один корень положительный — положение $\varphi = 0$ неустойчиво (седло).

\section*{5. Стратегия 1: стабилизация только угла}

\subsection*{5.1. Выбор управления}
\[
\ddot{u} = a\varphi + b\dot{\varphi}. \tag{3}
\]

\subsection*{5.2. Замкнутая система}
Подставляем (3) в (2):
\[
\ddot{\varphi} - \frac{g}{l}\varphi = -\frac{a\varphi + b\dot{\varphi}}{l},
\]
\[
\ddot{\varphi} + \frac{b}{l}\dot{\varphi} + \frac{a - g}{l}\varphi = 0. \tag{4}
\]

\subsection*{5.3. Анализ устойчивости}
Характеристическое уравнение:
\[
\lambda^2 + \frac{b}{l}\lambda + \frac{a - g}{l} = 0.
\]
Для асимптотической устойчивости необходимо:
\[
\frac{b}{l} > 0 \;\Rightarrow\; b > 0, \qquad \frac{a - g}{l} > 0 \;\Rightarrow\; a > g.
\]

\subsection*{5.4. Недостаток стратегии 1}
Управление не содержит $u$ и $\dot{u}$, поэтому точка опоры может не вернуться в ноль.

\section*{6. Стратегия 2: одновременная стабилизация угла и точки опоры}

\subsection*{6.1. Расширенное управление}
\[
\ddot{u} = -a\varphi - b\dot{\varphi} - cu - d\dot{u}, \quad a,b,c,d > 0. \tag{5}
\]

\subsection*{6.2. Система в переменных состояния}
Введем $x_1 = \varphi$, $x_2 = \dot{\varphi}$, $x_3 = u$, $x_4 = \dot{u}$:
\[
\begin{cases}
\dot{x}_1 = x_2, \\[2pt]
\dot{x}_2 = \dfrac{g}{l}x_1 - \dfrac{1}{l}x_4, \\[6pt]
\dot{x}_3 = x_4, \\[2pt]
\dot{x}_4 = -a x_1 - b x_2 - c x_3 - d x_4.
\end{cases}
\]

\subsection*{6.3. Характеристический полином}
\[
\lambda^4 + \left(d + \frac{b}{l}\right)\lambda^3 + \left(c + \frac{g}{l} + \frac{a}{l}\right)\lambda^2 + \frac{g}{l}d\lambda + \frac{g}{l}c = 0. \tag{6}
\]

\subsection*{6.4. Критерий Рауса-Гурвица}
Для полинома $\lambda^4 + A_3\lambda^3 + A_2\lambda^2 + A_1\lambda + A_0 = 0$ с $A_i > 0$ условие устойчивости:
\[
A_3 A_2 A_1 > A_1^2 + A_3^2 A_0.
\]

Для нашей системы:
\[
A_3 = d + \frac{b}{l},\quad A_2 = c + \frac{g}{l} + \frac{a}{l},\quad A_1 = \frac{g}{l}d,\quad A_0 = \frac{g}{l}c.
\]

Дополнительное условие:
\[
\left(d + \frac{b}{l}\right)\left(c + \frac{g}{l} + \frac{a}{l}\right)\left(\frac{g}{l}d\right) > 
\left(\frac{g}{l}d\right)^2 + \left(d + \frac{b}{l}\right)^2 \left(\frac{g}{l}c\right).
\]

\section*{7. Результаты экспериментов}

\subsection*{7.1. Автоматический режим}
При правильно подобранных коэффициентах наблюдается:
\[
\varphi(t) \to 0,\quad \dot{\varphi}(t) \to 0,\quad u(t) \to 0,\quad \dot{u}(t) \to 0.
\]

\subsection*{7.2. Влияние параметров}
\begin{itemize}
    \item Увеличение $a$ ускоряет возврат к вертикали.
    \item Увеличение $b$ подавляет колебания.
    \item Увеличение $c$ ускоряет возврат точки опоры.
    \item Увеличение $d$ демпфирует движение точки опоры.
\end{itemize}

\subsection*{7.3. Ручной режим}
При ручном управлении удается удерживать маятник только при медленном изменении.

\section*{8. Выводы}
\begin{enumerate}
    \item Неуправляемый перевернутый маятник неустойчив.
    \item Управление $\ddot{u} = a\varphi + b\dot{\varphi}$ с $a > g,\; b > 0$ делает верхнее положение устойчивым.
    \item Для стабилизации точки опоры нужно управление $\ddot{u} = -a\varphi - b\dot{\varphi} - cu - d\dot{u}$.
    \item Условия устойчивости задаются критерием Рауса-Гурвица.
\end{enumerate}

\end{document}